\section*{第 \ref{ch:g11:minerals} 章——矿物}

\begin{solution}{exo:g11:minerals:1}
\omterm{def:g11:minerals:cleavage}{解理}是矿物沿平整平面破裂的趋势，这些平面平行于
原子间键较弱的方向；\omterm{def:g11:minerals:cleavage}{断口}是
沿不规则表面的断裂，出现在不具弱面的矿物中。
石英没有\omterm{def:g11:minerals:cleavage}{解理}，\omterm{def:g11:minerals:cleavage}{断口}呈贝壳状，断面光滑弯曲，
像破碎的玻璃。云母有一组完全\omterm{def:g11:minerals:cleavage}{解理}，
易剥成\omterm{def:g11:minerals:thin-section}{薄片}。
\end{solution}

\begin{solution}{exo:g11:minerals:2}
微量其他元素即可改变石英的颜色，石英
可呈透明、白色、黑色、粉色、蓝色或紫色；橄榄石几乎总是
呈绿色。更可靠的判据有：\omterm{def:g11:minerals:cleavage}{解理}（或\omterm{def:g11:minerals:cleavage}{断口}）、硬度、\omterm{def:g11:minerals:lustre}{光泽}、
自由晶体的\omterm{def:g11:minerals:habit}{晶习}、密度。
\end{solution}

\begin{solution}{exo:g11:minerals:3}
（1）橄榄石。（2）黑云母。（3）正长石（长石类：两组\omterm{def:g11:minerals:cleavage}{解理}，硬度
6）。（4）石英（无\omterm{def:g11:minerals:cleavage}{解理}，硬度 7）。
\end{solution}

\begin{solution}{exo:g11:minerals:4}
（1）五个原子：中心一个硅原子，四个氧位于
四面体的顶角。（2）$+4 + 4 \times (-2) = -4$：基团为 \ce{SiO4^4-}。
\end{solution}

\begin{solution}{exo:g11:minerals:5}
橄榄石：孤立四面体。辉石：单链。角闪石：
双链。云母：层状。长石和石英：架状。
\end{solution}

\begin{solution}{exo:g11:minerals:6}
(1) $M(\ce{SiO2}) = 28.085 + 2 \times 15.999 = \qty{60.083}{g/mol}$;
氧：$31.998/60.083 \approx \qty{53.3}{\percent}$。
(2) $M(\ce{Mg2SiO4}) = 2 \times 24.305 + 28.085 + 4 \times 15.999 =
\qty{140.691}{g/mol}$；氧：$63.996/140.691 \approx \qty{45.5}{\percent}$。
（3）地壳中氧占 \qty{46}{\percent}，与之处于同一数量级：它几乎完全由富氧矿物
即\omterm{def:g11:minerals:silicate}{硅酸盐}构成。
\end{solution}

\begin{solution}{exo:g11:minerals:7}
白色颗粒具两组\omterm{def:g11:minerals:cleavage}{解理}，硬度介于 5.5 与 6.5 之间，属于
长石（6）。灰色颗粒硬度大于 6.5，\omterm{def:g11:minerals:cleavage}{断口}呈弯曲的贝壳状、无\omterm{def:g11:minerals:cleavage}{解理}：石英（7）。
%
\end{solution}

\begin{solution}{exo:g11:minerals:8}
（1）每个硅对应 $4 - 2/2 = 3$ 个氧：\ce{SiO3}。（2）\ce{(Si,Al)2O6} 含
两个硅（或铝）原子对应六个氧，每个硅对应三个氧，
相当于两个 \ce{SiO3}。
\end{solution}

\begin{solution}{exo:g11:minerals:9}
$\rho = 0.55 \times 2.69 + 0.35 \times 3.4 + 0.10 \times 3.32 = 1.480 + 1.190
+ 0.332 \approx \qty{3.00}{g/cm^3}$。立方体的体积为
\qty{1000}{cm^3}，故质量约为 \qty{3.00}{kg}。
\end{solution}

\begin{solution}{exo:g11:minerals:10}
(1) $0.75 \times 3.32 + 0.25 \times 3.4 = 2.49 + 0.85 = \qty{3.34}{g/cm^3}$.
（2）$(3.34 - 3.00)/3.00 \approx 0.11$：密度约高 \qty{11}{\percent}。
（3）\omterm{def:g11:ocean-ridges:oceanic-lithosphere}{洋壳}中的玄武岩和\omterm{def:g11:ocean-ridges:gabbro}{辉长岩}比
地幔中的\omterm{def:g11:ocean-ridges:peridotite}{橄榄岩}轻，它们浮在其上。
\end{solution}

\begin{solution}{exo:g11:minerals:11}
（1）斜长石 $168/400 = 0.42$，辉石 $0.33$，橄榄石 $0.15$，孔隙
$0.10$。（2）命题只对无孔隙的岩石成立；此处约十分之一
的体积是空的。（3）在固体部分，各组分比例为 $0.42/0.90
\approx 0.467$，$0.367$ 和 $0.167$：
$\rho_{\text{固体}} \approx 0.467 \times 2.69 + 0.367 \times 3.4 + 0.167
\times 3.32 \approx \qty{3.06}{g/cm^3}$。计入只增加体积而不增加质量的孔隙后，
岩石密度为 $0.42 \times 2.69 + 0.33 \times 3.4 +
0.15 \times 3.32 \approx \qty{2.75}{g/cm^3}$，即 $0.9 \times 3.06$。
\end{solution}

\begin{solution}{exo:g11:minerals:12}
（1）氧 $464/16.0 = 29.0$~mol；硅 $283/28.1 \approx 10.1$；铝
$84/27.0 \approx 3.1$；铁 $52/55.8 \approx 0.93$；钙 $46/40.1 \approx
1.15$；镁 $28/24.3 \approx 1.15$；钠 $23/23.0 = 1.0$；钾
$15/39.1 \approx 0.38$。（2）原子总计约 $46.8$~mol，其中
氧占 $29.0/46.8 \approx 0.62$：几乎每三个原子中就有两个是氧。
（3）体积较大的氧离子占原子的大部分，
在体积中所占比例更大：从原子尺度看，地壳是氧离子的堆积，
小的硅离子和金属离子填充其间的空隙。
\end{solution}

\begin{solution}{exo:g11:minerals:13}
（1）平均共享顶角数为 $(2 + 3)/2 = 2.5$，故每个硅对应
$4 - 2.5/2 = 2.75$ 个氧。（2）四个硅对应十一个氧：
\ce{Si4O11}。（3）\ce{Si7AlO22} 含八个四面体（七个硅、一个
铝）对应 22 个氧，即每个四面体 $22/8 = 2.75$ 个氧，或
两个 \ce{Si4O11}。两个 \ce{OH} 基团不属于四面体。
\end{solution}

\begin{solution}{exo:g11:minerals:14}
$0.20 \times 2.62 + 0.10 \times 2.56 + 0.50 \times 2.65 + 0.10 \times 3.09 +
0.10 \times 3.23 = 0.524 + 0.256 + 1.325 + 0.309 + 0.323 \approx
\qty{2.74}{g/cm^3}$，高于 \qty{2.7}{g/cm^3}；例中的花岗岩，
石英和钾长石更多，密度约为 \qty{2.59}{g/cm^3}，低于
\qty{2.6}{g/cm^3}。西部岩石密度更大，因其含有更多
斜长石、黑云母和角闪石，它们富含铁和镁；而高硅
花岗岩含更多石英和钾长石，这两种是最轻的
造岩矿物。
\end{solution}

\begin{solution}{exo:g11:minerals:15}
（1）上地壳：$66.6 \times 0.467 \approx \qty{31.1}{\percent}$ 硅；
全地壳：$60.6 \times 0.467 \approx \qty{28.3}{\percent}$。（2）上地壳
富含二氧化硅，应含更多石英和长石，而
富铁镁矿物更少；下地壳则相反。（3）石英
和长石是最轻的造岩矿物，因此上地壳应
比下地壳更轻。
\end{solution}

\begin{solution}{pb:g11:minerals:1}
\textbf{1.} 石英：无\omterm{def:g11:minerals:cleavage}{解理}，贝壳状\omterm{def:g11:minerals:cleavage}{断口}，玻璃\omterm{def:g11:minerals:lustre}{光泽}，硬度
高于指甲（7）。

\textbf{2.} 正长石，一种钾长石：两组\omterm{def:g11:minerals:cleavage}{解理}，
在莫氏硬度表上硬度为 6。

\textbf{3.} 斜长石（此处为更长石），
即含钠和钙的长石。

\textbf{4.} D 为黑云母，E 为白云母：两种云母。其\omterm{def:g11:minerals:silicate}{硅氧四面体}
构成片层，片层间的键较弱，
因此沿这一组完全\omterm{def:g11:minerals:cleavage}{解理}裂开。

\textbf{5.} 石英和两种长石（A、B、C）属架状结构；黑云母
（D）含铁和镁。

\textbf{6.} 石英 $150/500 = 0.30$；正长石 $0.35$；斜长石 $0.25$；
黑云母 $0.06$；白云母 $0.04$。

\textbf{7.} 若颗粒无序散布，落在某矿物上的点所占比例
估计该矿物在表面所占面积的比例，
而该比例等于其所占体积的比例。

\textbf{8.} $0.35 + 0.25 = 0.60$：长石共占 \qty{60}{\percent}。

\textbf{9.} 由 $1/500 = 0.002$，即 \qty{0.2}{\percent}。

\textbf{10.} 是：长石、石英和云母。

\textbf{11.} $\rho = \phi_1 \rho_1 + \phi_2 \rho_2 + \dots + \phi_k \rho_k$，
其中 $\phi_1 + \dots + \phi_k = 1$。

\textbf{12.} 石英 $0.30 \times 2.62 = 0.786$；正长石 $0.35 \times 2.56 =
0.896$；斜长石 $0.25 \times 2.65 = 0.6625$；黑云母 $0.06 \times 3.09
\approx 0.185$；白云母 $0.04 \times 2.82 \approx 0.113$（单位均为
\unit{g/cm^3}）。

\textbf{13.} $0.786 + 0.896 + 0.6625 + 0.185 + 0.113 \approx
\qty{2.64}{g/cm^3}$。

\textbf{14.} 正长石为 $0.896$，因其含量最多。黑云母和
白云母的密度分别为 \qty{3.09}{} 和 \qty{2.82}{g/cm^3}，
密度高于其他矿物，使岩石的密度高于
仅由石英和长石组成时的密度。

\textbf{15.} \qty{2640}{kg/m^3}：\qty{1}{m^3} 花岗岩的质量
约为 \qty{2.64}{t}。

\textbf{16.} $100 \times 2.643 \approx \qty{264.3}{g}$。

\textbf{17.} 岩块排开 \qty{100}{cm^3} 的水，水的质量为
\qty{100}{g}：天平示数为 $264.3 - 100 = \qty{164.3}{g}$。

\textbf{18.} $264.3/(264.3 - 164.3) = 264.3/100 \approx \qty{2.64}{g/cm^3}$。

\textbf{19.} 介于两者之间：比平均值 \qty{2.69}{g/cm^3} 轻，
但高于二氧化硅和钾含量最高的花岗岩的 \qty{2.6}{g/cm^3}。
增加斜长石以及黑云母或角闪石、减少石英和钾长石，
会使花岗岩的密度更大。

\textbf{20.} 该花岗岩的密度约为
\textbf{\qty{2.64}{g/cm^3}}。
\end{solution}
